\section{Segmentación: PCA + KMeans}
\label{sec:segmentacion}

La segmentación parte de un panel por cliente con ocho \emph{features}:
\texttt{net\_revenue}, \texttt{margin\_rate}, \texttt{frequency},
\texttt{retention\_rate}, \texttt{aov}, \texttt{recency\_days},
\texttt{return\_rate} y \texttt{cltv}. El pipeline es la composición
canónica $X\to Z=\text{StandardScaler}(X)\to Y=\text{PCA}_{n=2}(Z)\to
\text{KMeans}_{k=3}(Y)$, encapsulada en
\texttt{src/analytics/clustering.py}.

\subsection*{Fundamento matemático}

Sea $S=\tfrac{1}{n-1}Z^{\top}Z$ la covarianza muestral. Las direcciones
óptimas del PCA verifican

\begin{align}
  w_i \;=\; \argmax_{\|w\|=1}\; w^{\top}\,S\,w,
\end{align}

cuya solución por las condiciones KKT da $S\,w_i = \lambda_i\,w_i$: los
\emph{vectores propios} de $S$ ordenados por \emph{valor propio}
descendente. Las dos primeras componentes explican $57{,}1\%$ y
$23{,}2\%$ de la varianza ($\Sigma \approx 80{,}3\%$), valores leídos del
atributo \texttt{explained\_variance\_ratio\_} de
\texttt{models/pca.joblib}.

El KMeans minimiza la inercia intra-cluster

\begin{align}
  J(C_1,\dots,C_k)
  \;=\;
  \sum_{i=1}^{k}\;\sum_{y\in C_i}\;\bigl\|y - \mu_i\bigr\|_2^{2},
  \qquad
  \mu_i \;=\; \tfrac{1}{|C_i|}\sum_{y\in C_i}y,
  \label{eq:kmeans-inertia}
\end{align}

mediante el algoritmo de Lloyd, entrenado sobre el espacio reducido
$Y\in\mathbb{R}^{n\times 2}$. La curva del codo (figura~\ref{fig:elbow})
cae de $13\,360$ ($k=2$) a $2\,934$ ($k=3$) y desciende suavemente desde
ahí; se fija $k^{\star} = 3$.

\begin{figure}[H]
  \centering
  \includegraphics[width=0.65\linewidth]{elbow.png}
  \caption{Inercia $J^{\star}(k)$ frente a $k\in\{2,\dots,8\}$. El codo en
  $k^{\star}=3$ justifica la elección final.}
  \label{fig:elbow}
\end{figure}

\subsection*{Resultado: perfiles de cluster}

\begin{figure}[H]
  \centering
  \includegraphics[width=0.78\linewidth]{pca_clusters.png}
  \caption{Proyección PC1--PC2 de los $5\,750$ clientes coloreada por
  cluster KMeans.}
  \label{fig:pca-clusters}
\end{figure}

\begin{table}[H]
\centering
\caption{Perfil mediano por cluster, leído de
\texttt{models/customer\_segments.parquet}.}
\label{tab:cluster-perfil}
\renewcommand{\arraystretch}{1.25}
\resizebox{\linewidth}{!}{%
\begin{tabular}{l r r r r r r r r r}
\toprule
\textbf{Cluster} & \textbf{$n$} &
\texttt{net\_rev.} & \texttt{margin} & \texttt{freq.} &
\texttt{ret.} & \texttt{aov} & \texttt{recency} &
\texttt{ret\_rate} & \textbf{CLTV} \\
\midrule
0 (\emph{ocasional})   & $4\,576$ & $129{,}90$    & $0{,}400$ & $1$  & $0{,}014$ & $129{,}90$ & $854$ & $0{,}000$ & $0{,}72$        \\
1 (\emph{alto valor})  & $740$    & $11\,461{,}53$ & $0{,}388$ & $20$ & $0{,}250$ & $572{,}24$ & $82$  & $0{,}022$ & $21\,191{,}30$ \\
2 (\emph{devolutivo})  & $434$    & $0{,}00$       & $0{,}000$ & $1$  & $0{,}014$ & $0{,}00$   & $856$ & $1{,}000$ & $0{,}00$        \\
\bottomrule
\end{tabular}}
\end{table}

El \emph{cluster 0} ($\approx 80\%$ de la base) recoge los clientes
\emph{one-shot}. El \emph{cluster 1} ($\approx 13\%$) son recurrentes con
mediana de $20$ tickets, AOV de $572$\,\texteuro\ y retención del $25\%$,
y concentran el grueso del CLTV agregado. El \emph{cluster 2}
($\approx 7{,}5\%$) son clientes devolutivos con tasa de devolución
mediana del $100\%$ — el grupo separado por PC2, y ruido desde el punto
de vista de valor.
